\subsection{RQ1: Does process-level organization help?}
\label{sec:results-effectiveness}

Table~\ref{tab:pairwise-results} reports the primary pairwise outcomes. Under same-family judging, \condFull\ is preferred to all three baseline organizations: its pooled mean win rate across replications is \PooledFourOneRateMean\ against \condSinglePass, \PooledFourTwoRateMean\ against \condStaged, and \PooledFourThreeRateMean\ against \condTaskPlanning. The direction is stable across benchmarks and replications: \condFull\ is favored in every benchmark-by-replication cell against \condSinglePass\ and \condTaskPlanning, and in all but one cell against \condStaged. After Holm correction, \PooledFourOneHolmCellsTotal, \PooledFourTwoHolmCellsTotal, and \PooledFourThreeHolmCellsTotal\ of the nine cells survive for the three contrasts, respectively. Replicate-level intervals are shown in Appendix~\ref{app:pairwise-detail}.

\begin{table*}[t]
\centering
\small
\setlength{\tabcolsep}{5pt}
\begin{tabular}{lrrrr}
\toprule
\textbf{Contrast}
& \textbf{WritingBench}
& \textbf{HelloBench}
& \textbf{DoLoMiTes}
& \textbf{Across benchmarks} \\
\midrule
\condAgenticCogWriter\ vs \condSinglePass
& \WritingFourOneRateMean
& \HelloFourOneRateMean
& \DoLoFourOneRateMean
& \PooledFourOneRateMean \\
\condAgenticCogWriter\ vs \condStaged~\cite{yang2022re3,wan2025cognitive}
& \WritingFourTwoRateMean
& \HelloFourTwoRateMean
& \DoLoFourTwoRateMean
& \PooledFourTwoRateMean \\
\condAgenticCogWriter\ vs \condTaskPlanning~\cite{xiong2025beyond}
& \WritingFourThreeRateMean
& \HelloFourThreeRateMean
& \DoLoFourThreeRateMean
& \PooledFourThreeRateMean \\
\bottomrule
\end{tabular}
\caption{Primary prompt-collapsed pairwise outcomes against the controlled baseline organizations. Prior-work citations indicate the organizational precedent represented by each structured baseline; all conditions are our controlled re-implementations. Benchmark columns report Agentic CogWriter's arithmetic mean of replicate-specific win rates among non-tied comparisons. For the final column, wins, losses, and ties are first pooled across benchmarks within each replication, and the resulting replicate-level win rates are then averaged. These means are descriptive summaries; confirmatory sign tests, replicate-specific Wilson intervals, and Holm decisions are reported in Appendix~\ref{app:pairwise-detail}.}
\label{tab:pairwise-results}
\end{table*}


\subsection{RQ2: Which control mechanisms matter?}
\label{sec:results-mechanism}

The ablations separate the full architecture from two of its dynamic control commitments. Removing the explicit goal network yields a bounded null: \condFull\ wins \PooledFourFiveRateMean\ of non-tied comparisons against \condNoGoals, and no benchmark cell survives Holm correction. Fixing the process order is similarly close to \condFull, with no surviving benchmark cell. These results do not establish equality, but they provide no evidence that either explicit goal bookkeeping or adaptive process ordering accounts for the main advantage.

The traces are consistent with this product-level result. The unrestricted \texttt{Monitor} follows the canonical \texttt{Planning}$\rightarrow$\texttt{Translating}$\rightarrow$\texttt{Reviewing} cycle in \AfourCycleCompliantRate\ of completed runs, returns from \texttt{Reviewing} to \texttt{Planning} in only \AfourReviewingToPlanningRate, and accepts goal regeneration in only \AfourRegenerationRuns\ of \AfourRunCount\ runs. Collapsing the process roles into \condSingleWriter\ causes a much larger drop: \condFull\ is preferred in \PooledFourSevenRateMean\ of non-tied comparisons. That contrast is not isolating, however, because it also changes recursion granularity and output length. Detailed process counts and sequence distributions appear in Appendix~\ref{app:process-behavior}.

\subsection{Robustness}
\label{sec:results-robustness}

The main pattern is not straightforwardly explained by output length or inference expenditure. The same-family judge prefers the longer output in \LengthLongerRate\ of non-tied pairs, but \condFull\ remains favored over \condSinglePass\ within the tested output-length matching bands. At a compute ratio within 1.25, \condFull\ remains preferred to \condStaged\ at \ComputeFullStagedWithinOneTwentyFiveRate\ and to \condTaskPlanning\ at \ComputeFullTaskPlanningWithinOneTwentyFiveRate; there are no \condFull--\condSinglePass\ pairs this closely matched in compute. These realized-compute analyses are descriptive rather than causal because compute is itself an outcome of execution.

A cross-family judge preserves the pooled direction for the three checked contrasts, including \CrossFamilyFullSinglePassWinRate\ for \condFull\ versus \condSinglePass\ and \CrossFamilyFullTaskPlanningWinRate\ versus \condTaskPlanning, but commits on only \CrossFamilyPooledCommitRate\ of eligible pairs and reverses \condFull\ versus \condSinglePass\ on WritingBench. We therefore treat it as a robustness check rather than replacement confirmatory evidence. Native and pointwise scores, length matching, compute strata, and cross-family counts are reported in Appendix~\ref{app:secondary-scores}--\ref{app:cross-family}.
